\textbf{Phần 1. Trắc nghiệm lựa chọn}
\setcounterpageref{ex}{0}

\begin{ex} %[3D6N39]
Muốn so sánh số lớn gấp mấy lần số bé, ta lấy
\choice
{\True Số lớn chia cho số bé}
{Số lớn cộng số bé}
{Số lớn trừ số bé}
{Số lớn nhân số bé}
\loigiai{
Muốn tìm số lớn gấp mấy lần số bé: Số lớn : Số bé.\\
\textbf{Đáp án của câu hỏi là A}
}
\end{ex}

\begin{ex} %[3D6N39]
Số 24 gấp mấy lần số 6?
\choice
{$3$ lần}
{$5$ lần}
{$6$ lần}
{\True $4$ lần}
\loigiai{
$24 : 6 = 4$ lần.\\
\textbf{Đáp án của câu hỏi là D}
}
\end{ex}

\begin{ex} %[3D6H39]
Số 45 gấp mấy lần số 5?
\choice
{$8$ lần}
{$7$ lần}
{\True $9$ lần}
{$10$ lần}
\loigiai{
$45 : 5 = 9$ lần.\\
\textbf{Đáp án của câu hỏi là C}
}
\end{ex}

\begin{ex} %[3D6H39]
Đoạn thẳng $AB$ dài $32\text{ cm}$, đoạn thẳng $CD$ dài $8\text{ cm}$. Đoạn thẳng $AB$ dài gấp mấy lần đoạn thẳng $CD$?
\choice
{$3$ lần}
{\True $4$ lần}
{$5$ lần}
{$6$ lần}
\loigiai{
$32 : 8 = 4$ lần.\\
\textbf{Đáp án của câu hỏi là B}
}
\end{ex}

\begin{ex} %[3D6H39]
Bố 36 tuổi, con 6 tuổi. Tuổi bố gấp mấy lần tuổi con?
\choice
{$5$ lần}
{$7$ lần}
{$8$ lần}
{\True $6$ lần}
\loigiai{
$36 : 6 = 6$ lần.\\
\textbf{Đáp án của câu hỏi là D}
}
\end{ex}

\begin{ex} %[3D6H39]
Mẹ có 56 quả táo, bà có 8 quả táo. Số táo của mẹ gấp mấy lần số táo của bà?
\choice
{$6$ lần}
{$8$ lần}
{\True $7$ lần}
{$9$ lần}
\loigiai{
$56 : 8 = 7$ lần.\\
\textbf{Đáp án của câu hỏi là C}
}
\end{ex}

\begin{ex} %[3D6V39]
Thùng thứ nhất có $48\text{ l}$ dầu, thùng thứ hai có $6\text{ l}$ dầu. Hỏi số dầu ở thùng thứ nhất gấp mấy lần số dầu ở thùng thứ hai?
\choice
{$7$ lần}
{$9$ lần}
{$6$ lần}
{\True $8$ lần}
\loigiai{
$48 : 6 = 8$ lần.\\
\textbf{Đáp án của câu hỏi là D}
}
\end{ex}

\begin{ex} %[3D6V39]
Một hình chữ nhật có chiều dài $28\text{ cm}$, chiều rộng $7\text{ cm}$. Chiều dài gấp mấy lần chiều rộng?
\choice
{\True $4$ lần}
{$3$ lần}
{$5$ lần}
{$6$ lần}
\loigiai{
$28 : 7 = 4$ lần.\\
\textbf{Đáp án của câu hỏi là A}
}
\end{ex}

\begin{ex} %[3D6V39]
Đoạn dây đỏ dài $72\text{ cm}$, đoạn dây xanh dài $9\text{ cm}$. Đoạn dây đỏ dài gấp đoạn dây xanh số lần là
\choice
{$7$ lần}
{\True $8$ lần}
{$9$ lần}
{$10$ lần}
\loigiai{
$72 : 9 = 8$ lần.\\
\textbf{Đáp án của câu hỏi là B}
}
\end{ex}

\begin{ex} %[3D6V39]
Lớp 3A có 32 học sinh, trong đó có 8 học sinh xuất sắc. Tổng số học sinh cả lớp gấp số học sinh xuất sắc số lần là
\choice
{\True $4$ lần}
{$3$ lần}
{$5$ lần}
{$6$ lần}
\loigiai{
$32 : 8 = 4$ lần.\\
\textbf{Đáp án của câu hỏi là A}
}
\end{ex}

\begin{ex} %[3D6V39]
Một bao gạo nặng $60\text{ kg}$, một túi gạo nặng $5\text{ kg}$. Bao gạo nặng gấp túi gạo số lần là
\choice
{$10$ lần}
{$14$ lần}
{\True $12$ lần}
{$15$ lần}
\loigiai{
$60 : 5 = 12$ lần.\\
\textbf{Đáp án của câu hỏi là C}
}
\end{ex}

\begin{ex} %[3D6C39]
Năm nay mẹ 35 tuổi, con 7 tuổi. Hỏi sau 7 năm nữa tuổi mẹ gấp mấy lần tuổi con?
\choice
{$2$ lần}
{\True $3$ lần}
{$4$ lần}
{$5$ lần}
\loigiai{
Sau 7 năm: Mẹ $35 + 7 = 42$ tuổi; Con $7 + 7 = 14$ tuổi.\\
Số lần: $42 : 14 = 3$ lần.\\
\textbf{Đáp án của câu hỏi là B}
}
\end{ex}

\textbf{Phần 2. Trắc nghiệm đúng sai}
\setcounterpageref{ex}{0}

\begin{ex} %[3D6H39]
Xét tính đúng sai về so sánh số lớn gấp mấy lần số bé:
\choiceTF
{\True Số 36 gấp 4 lần số 9 ($36 : 9 = 4$).}
{\True Số 40 gấp 8 lần số 5 ($40 : 5 = 8$).}
{Số 30 gấp 6 lần số 5 là phép so sánh hơn kém số đơn vị.}
{\True Số 30 hơn số 5 là 25 đơn vị, và gấp 6 lần số 5.}
\loigiai{
\begin{itemize}
    \item $36 : 9 = 4$: Đúng.
    \item $40 : 5 = 8$: Đúng.
    \item Gấp mấy lần là so sánh tỉ lệ (phép chia), không phải hơn kém đơn vị. Ý 3 Sai, Ý 4 Đúng.
\end{itemize}
\begin{itemize}
    \item a) \textbf{Đúng}.
    \item b) \textbf{Đúng}.
    \item c) \textbf{Sai}.
    \item d) \textbf{Đúng}.
\end{itemize}
}
\end{ex}

\begin{ex} %[3D6V39]
Đoạn thẳng $AB = 48\text{ cm}$, đoạn thẳng $CD = 6\text{ cm}$. Xét tính đúng sai:
\choiceTF
{\True Đoạn $AB$ dài gấp 8 lần đoạn $CD$.}
{\True Đoạn $CD$ bằng $\dfrac{1}{8}$ đoạn $AB$.}
{Đoạn $AB$ dài hơn đoạn $CD$ là 8 cm.}
{\True Đoạn $AB$ dài hơn đoạn $CD$ là 42 cm ($48 - 6 = 42$).}
\loigiai{
\begin{itemize}
    \item $48 : 6 = 8$ lần. Đúng.
    \item $CD = \dfrac{1}{8} AB$. Đúng.
    \item $AB$ hơn $CD$: $48 - 6 = 42\text{ cm} \neq 8\text{ cm}$. Ý 3 Sai, Ý 4 Đúng.
\end{itemize}
\begin{itemize}
    \item a) \textbf{Đúng}.
    \item b) \textbf{Đúng}.
    \item c) \textbf{Sai}.
    \item d) \textbf{Đúng}.
\end{itemize}
}
\end{ex}

\textbf{Phần 3. Trả lời ngắn}
\setcounterpageref{ex}{0}

\begin{ex} %[3D6H39]
Số 54 gấp mấy lần số 6?
\shortans{9}
\loigiai{
\textbf{Đáp số:} $9$\\
$54 : 6 = 9$ lần.
}
\end{ex}

\begin{ex} %[3D6H39]
Số 63 gấp mấy lần số 7?
\shortans{9}
\loigiai{
\textbf{Đáp số:} $9$\\
$63 : 7 = 9$ lần.
}
\end{ex}

\begin{ex} %[3D6V39]
Bà 64 tuổi, cháu 8 tuổi. Hỏi tuổi bà gấp mấy lần tuổi cháu?
\shortans{8}
\loigiai{
\textbf{Đáp số:} $8$\\
$64 : 8 = 8$ lần.
}
\end{ex}

\begin{ex} %[3D6V39]
Một tấm vải hoa dài $45\text{ m}$, tấm vải trắng dài $5\text{ m}$. Hỏi tấm vải hoa dài gấp mấy lần tấm vải trắng?
\shortans{9}
\loigiai{
\textbf{Đáp số:} $9$\\
$45 : 5 = 9$ lần.
}
\end{ex}

\begin{ex} %[3D6V39]
Một hình chữ nhật có chu vi $32\text{ cm}$, một hình vuông có chu vi $8\text{ cm}$. Chu vi hình chữ nhật gấp mấy lần chu vi hình vuông?
\shortans{4}
\loigiai{
\textbf{Đáp số:} $4$\\
$32 : 8 = 4$ lần.
}
\end{ex}

\begin{ex} %[3D6C39]
Có 42 quả cam và 6 quả bưởi. Cần thêm bao nhiêu quả bưởi nữa để số cam chỉ còn gấp 3 lần số bưởi?
\shortans{8}
\loigiai{
\textbf{Đáp số:} $8$\\
Để số cam gấp 3 lần số bưởi thì số bưởi cần có là: $42 : 3 = 14$ quả.\\
Số bưởi cần thêm là: $14 - 6 = 8$ quả.
}
\end{ex}
